% IEEEtran selects the legacy `ptm` family while the class is loading.  XeTeX
% uses TU encoding, so provide silent bootstrap shapes until fontspec installs
% the embedded OpenType family below.
\DeclareFontFamily{TU}{ptm}{}
\DeclareFontShape{TU}{ptm}{m}{n}{<->ssub*lmr/m/n}{}
\DeclareFontShape{TU}{ptm}{bx}{n}{<->ssub*lmr/bx/n}{}
\DeclareFontShape{TU}{ptm}{m}{it}{<->ssub*lmr/m/it}{}
\DeclareFontShape{TU}{ptm}{bx}{it}{<->ssub*lmr/bx/it}{}
\documentclass[conference]{IEEEtran}

\usepackage{iftex}
\ifXeTeX\else
  \PackageError{paper-build}{This paper requires XeLaTeX}{Run the PowerShell build script}
\fi

\usepackage{xeCJK}
\setmainfont{Times New Roman}
\setCJKmainfont[
  BoldFont=FandolSong-Bold,
  ItalicFont=FandolKai-Regular,
  BoldItalicFont=FandolHei-Bold,
  SmallCapsFont=FandolSong-Regular
]{FandolSong-Regular}
\setCJKsansfont[BoldFont=FandolHei-Bold]{FandolHei-Regular}
\setCJKmonofont{FandolFang-Regular}
\usepackage{amsmath,amssymb,mathtools}
\usepackage{booktabs}
\usepackage{array}
\usepackage{tabularx}
\usepackage{graphicx}
\usepackage{xcolor}
\usepackage{url}
\usepackage{tikz}
\usepackage{pgfplots}
\pgfplotsset{compat=1.18}
\usetikzlibrary{arrows.meta,calc,fit,matrix,positioning,shapes.geometric}
\usepackage[hidelinks]{hyperref}
\hypersetup{
  pdftitle={从多方向到无显式方向：任意秩拒答消融的形式化、适配与失败实验分析},
  pdfauthor={匿名作者},
  pdfsubject={拒答消融的历史形式化、本地适配与验证集失败实验},
  pdfkeywords={拒答消融, 源码形式化, 局部邻域映射, 多目标优化}
}

\newcolumntype{Y}{>{\raggedright\arraybackslash}X}
\newcolumntype{C}{>{\centering\arraybackslash}X}
\newcommand{\code}[1]{\texttt{#1}}
\newcommand{\method}{ARA}
\newcommand{\sommethod}{SOM 多方向消融}
\urlstyle{same}

\title{从多方向到无显式方向：任意秩拒答消融的形式化、适配与失败实验分析}

\author{\IEEEauthorblockN{匿名作者（投稿占位）}
\IEEEauthorblockA{匿名单位（投稿占位）\\
匿名邮箱（投稿占位）}}

\begin{document}

\maketitle

\begin{abstract}
\label{sec:abstract}
语言模型拒答行为的权重干预需要同时回答两个问题：局部表示如何改变，改变后能否满足预设的行为与分布约束。本文以 Heretic 方向消融为起点，在注明上游任意秩消融来源的基础上，研究本地校准式实现及其后续适配。单位置版本冻结基座输入输出，通过秩上界为 128 的增量、尺度归一化的平滑邻居距离、保持与间隔推远目标及因子规范化优化局部模块；轨迹版本进一步实现逐步对齐采集、部署增益和隔离审计协议。本文离线重放一次 Qwen3.8-27B 归档搜索的全部 120 个候选，全部状态为完成，但无候选通过验收门槛。验证集最低拒答关键词命中率为 54/100，对应首 token KL 为 0.089975 nats；该候选并非独立测试结果或已验收模型。所有候选满足 KL 数值门槛，却均未满足关键词门槛，独立审计及重载分数无记录。结果说明，这次单位置搜索产生了局部行为变化，但未达到既定干预目标；低首 token 漂移不能替代完整能力评价。轨迹版本仅报告已实现机制，本归档没有其效果数据。本文据此给出参数化、数值实现与证据边界一致的方法分析，为后续同协议对照提供可复查基础。

\noindent\textbf{关键词：}机制可解释性；拒答行为；任意秩消融；低秩适配；实验复核
\end{abstract}

\section{引言}
\label{sec:introduction}

拒答行为的可干预性为理解语言模型内部表示提供了一个具体研究问题：当部分权重沿某种表示方向改变时，哪些输出行为随之变化，哪些性质仍保持稳定。Arditi 等在所研究的模型与设置中发现，拒答行为可以由一个方向介导，并展示了相应干预\cite{arditi2024}。这一结果支持对拒答相关表示开展机制分析，但不能推导出所有模型、所有语言或所有请求都共享同一条充分的拒答方向。对一个新的模型实例，干预空间与评价协议仍需分别检验。

Heretic 将方向消融组织为自动化搜索流程：从两组提示的残差统计中构造方向，再搜索层位置、组件权重和干预范围，以拒答相关指标与输出分布漂移共同评价候选\cite{heretic}。这种参数化容易解释，也使候选之间能够保持统一接口。其限制在于，局部更新受到所选方向及权重核约束；即便搜索充分，也只覆盖由该表示假设定义的更新族。因此，扩大更新空间是合理的方法动机，但“更大的空间”本身并不意味着更好的实测结果。

任意秩消融（Arbitrary-Rank Ablation，ARA）在上游实现中转向模块输入输出，利用保持、拉近与推远目标求解映射变化\cite{upstreamara}。本文沿用这一思想及名称，把研究对象限定为本地校准式 ARA（Calibrated Arbitrary-Rank Ablation，下文简称本地 ARA；配置字段仍称 CARA）。本地实现通过有限秩因子表达增量，并对距离尺度、推远目标和部署因子进行处理。这里的贡献是对既有思想的适配、明确化和证据分析，不把 ARA 名称、三项基本目标或 LoRA 路径归为本项目首创。

本文区分已经运行的单位置版本 point-v1 与当前实现的轨迹版本 trajectory-v2。前者只采集提示末端预测首 token 的模块状态，后者使用冻结基座短续写的逐步状态，并增加部署增益、连续评分及导出约束。两者不仅采集位置不同，校准规模、验证切分、关键词集合和搜索协议也有差异。因此，当前源码只能说明后续机制已经实现，不能用来重新命名历史结果，更不能把多个协议下的数字汇入同一条效果结论。

实验结论是一个明确的未达标结果。唯一量化来源为 \code{docs/logs/ara-v1/} 的归档：一次搜索包含 120 个候选，全部完成，但无候选通过验收。按验证集关键词命中率优先排序，trial 66 的结果为 54/100，对应首 token KL 为 0.089975；机器验收没有选中该候选。最低 KL 来自另一个候选，二者不能拼接。未修改基座可作为同次运行对照，Heretic 方向消融则没有同协议实测数据。这样的证据支持分析本次搜索的失败位置，不支持判断本地 ARA 优于 Heretic。

本文的工作包含三个相互衔接的部分。方法部分把本地单位置目标写为与源码一致的数学形式，解释有限秩边界、非凸优化和规范化条件。适配部分说明轨迹采集、模块识别与验收状态之间的关系，明确哪些变化尚未获得效果证据。实验部分从原始 journal 重建候选分布，区分进程完成、数值门槛和模型验收，并讨论代理指标与选择偏差。论文的价值在于使方法描述和实验结论具有同一证据边界，而不是通过放宽验收条件把负结果改写为成功。

\section{方法背景与相关研究}
\label{sec:background}

\subsection{从拒答相关方向到权重更新}

方向消融先把行为差异转为表示差异。令第 $\ell$ 层的 good 与 bad 残差均值分别为 $\mu_{g,\ell}$ 和 $\mu_{b,\ell}$；在非零差向量条件下，单位方向为
\begin{equation}
 d_\ell=\frac{\mu_{b,\ell}-\mu_{g,\ell}}
 {\|\mu_{b,\ell}-\mu_{g,\ell}\|_2}.
 \label{eq:direction}
\end{equation}
good 与 bad 表示校准提示所属的数据组，并不意味着模块中的每个坐标都具有相应语义。当前 \code{workflow.py} 的方向准备函数计算该均值差；启用可选正交化时，还会去除 good 均值方向上的投影并重新归一化。方向因而是由数据和设置共同决定的统计对象。

Heretic 可逐层采用 $d_\ell$，也可在相邻层方向间线性插值后重新归一化，得到共享干预方向。组件权重核在选定中心达到最大值，随层距线性下降，到区间外跳过该层。令模块权重为 $W_0\in\mathbb{R}^{d_{\rm out}\times d_{\rm in}}$，输出空间单位方向为 $d$，组件权重为 $\lambda$。仅在不做行归一化的 \code{NONE} 模式下，局部更新可写为
\begin{equation}
 W'=W_0-\lambda dd^\top W_0.
 \label{eq:direction-update}
\end{equation}
这对应秩至多为一的增量。\code{PRE} 先在归一化权重上构造更新，再按原行范数缩放；\code{FULL} 进一步规范化更新后的行并近似分解增量。因此式~\eqref{eq:direction-update} 不能替代全部配置的实现定义。

这类干预的解释范围由构造方式决定。一个方向能够改变已测提示上的输出，说明该方向与相关行为具有可干预联系；要证明它是跨提示、跨模型的充分机制，还需要额外的干预与对照。本文借用方向方法作为动机和实现参照，不把其局限预设为本次 ARA 搜索失败或成功的原因。

\subsection{任意秩消融与低秩参数化}

上游 ARA 的关键变化是把优化对象移到模块映射。固定提交中的实现捕获 good/bad 输入输出，在保留 good 输出的同时，把更新后的 bad 输出拉向 good 参考，并以负距离项推离原 bad 参考。该版本同时存在全矩阵求解与 LoRA 求解两条路径\cite{upstreamara}。本文的方法对照据此区分目标形式和数值处理，不使用“上游只有全矩阵”的归属判断。

低秩适配（Low-Rank Adaptation，LoRA）用矩阵乘积表示冻结基座上的权重增量\cite{hu2022lora}。对于输入、输出维度较大的投影，因子参数量为 $r(d_{\rm in}+d_{\rm out})$，小于完整矩阵参数量的条件是 $r<d_{\rm in}d_{\rm out}/(d_{\rm in}+d_{\rm out})$。这种参数化与是否采用任务标签、何种局部损失是不同问题。本地 ARA 使用 LoRA 承载更新，但局部目标来自校准输出几何关系；引用 LoRA 说明参数化背景，并不构成创新性证明。

\subsection{行为代理与有效性证据}

拒答关键词和输出分布衡量不同对象。关键词命中率来自模型生成文本中的字符串匹配，首 token KL 来自无害提示上两个词表分布之间的差异。前者能呈现某些拒答措辞的频率，后者能描述输出开端的漂移；二者均不足以单独判定整段回答的正确性、有害性或任务能力。本文保留这种区别，避免把指标名称直接替换成未经测量的语义概念。

同样，局部优化、候选搜索与验收需要不同证据。损失下降发生在固定模块输入输出上，验证分数发生在修改后的全模型上，独立审计则应发生在候选锁定之后。优化代码存在不能替代候选实测，候选分数达标也不能替代独立审计和导出重载。后文按这三个层次组织公式、协议和结果，使读者能够定位每项结论依赖的观测。

\section{上游 ARA 原型的历史形式化}
\label{sec:method}

本节仅分析固定于 \code{edc3b123456c} 的上游原型。下述硬 $k$NN、全矩阵行归一化、五次 L-BFGS 调用、低秩路径仅重置 $B$ 与分段外层评分均为历史实现事实，不描述第~\ref{sec:ara-v1-results} 节的本地实验。本地 point-v1 的不同公式与控制流见第~\ref{sec:local-adaptation} 节。

\subsection{双层流程与问题定义}

本节把源实现还原为一个“冻结表示、内层改写、外层选择”的双层过程，而不把代码行为扩展为未经验证的训练算法。给定一组需要保持正常响应的提示 $\mathcal P_g$ 和一组需要降低拒答倾向的目标提示 $\mathcal P_b$，方法先在原模型上缓存所选模块的输入与输出，再独立修改各模块的线性映射，最后以模型级代理指标选择层区间和损失超参数。图~\ref{fig:ara-pipeline} 给出数据流：提示只在缓存阶段经过原模型，内层 L-BFGS 随后反复读取固定张量，不在每个闭包中重新运行完整 Transformer。该分离使参数更新可在模块尺度解释，也意味着内层目标只描述缓存点附近的行为，不能自动推出对任意提示的全局保证。源实现及其固定修订版本是本节历史算法事实的依据~\cite{ara2026source}。

设模型共有 $L$ 层，外层给出半开层区间 $[\ell_s,\ell_e)$，并在其中选择注意力输出投影和 MLP 下投影。对任一目标线性模块 $m$，记输入维为 $p$、输出维为 $q$，原权重为 $W_0\in\mathbb R^{q\times p}$。本节省略模块下标；实际过程按模块顺序重复。外层变量写为
\begin{equation}
  \theta=(\ell_s,\ell_e,\lambda_g,\lambda_b,\beta,k),
\end{equation}
其中 $\lambda_g$ 控制正常行为保持，$\lambda_b$ 控制目标提示的局部重映射，$\beta$ 控制远离原目标输出邻域的校正幅度。区间端点与连续参数的实际搜索域列于表~\ref{tab:source-parameters}；这些数值是源码默认值和搜索空间，不是本文重新调参所得的实验设置。

一次候选配置的执行顺序可以进一步展开为四步。首先，从未修改的基线模型收集所选层中全部目标模块的末词元记录，并按提示类别形成两个缓存字典。其次，对字典中的每个模块建立自己的优化变量，使用同一组 $\theta$ 最小化模块级目标；一个模块的损失不读取其他模块的权重，也不在上一模块改写后重新采集其输入。再次，将所有模块改写共同装入模型，在独立评估提示上计算外层两个目标。最后，外层搜索器记录该点并提出下一候选，模型快速恢复到基线映射后进入下一试次。于是，模块层面的拟合在冻结样本上可分解，模型层面的效果却由多个改写的组合决定；外层评分承担了检查这种组合结果的职责。

这一流程还区分了三种容易混淆的“方向”。参数优化当然产生梯度方向，L-BFGS 也保存历史搜索方向，但方法没有在表示空间中预先定义一个对所有目标样本共享的语义消融向量。真正决定每个目标输出移动参照的是它在 $Y_g$ 与 $Y_b$ 中的当前近邻，参照集合可因样本、模块和迭代而异。本文所称“无显式方向”专指没有固定的全局表示方向，而不是声称数值优化不使用方向信息，也不是声称任意修改都与子空间结构无关。

\begin{figure*}[t]
  \centering
  \resizebox{\textwidth}{!}{%
  \begin{tikzpicture}[
    font=\small,
    box/.style={draw, rounded corners=2pt, align=center, minimum height=11mm,
      minimum width=25mm, fill=blue!5, inner sep=4pt},
    cache/.style={box, fill=cyan!8},
    opt/.style={box, fill=orange!10},
    map/.style={box, fill=green!9},
    score/.style={box, fill=purple!8},
    arr/.style={-{Latex[length=2.2mm]}, thick},
    feedback/.style={-{Latex[length=2.2mm]}, thick, dashed, color=purple!70!black}
  ]
    \node[box] (prompt) {保持提示集 $\mathcal P_g$\\目标提示集 $\mathcal P_b$};
    \node[box, right=7mm of prompt] (hook) {目标层前向钩子\\仅取末词元};
    \node[cache, right=7mm of hook] (iocache) {冻结 I/O 缓存\\$X_g,Y_g,X_b,Y_b$};
    \node[opt, right=7mm of iocache] (loss) {$k$ 近邻三项损失\\$\mathcal L_{\rm total}$};
    \node[opt, right=7mm of loss] (lbfgs) {L-BFGS 内层优化\\强 Wolfe 线搜索};

    \node[map, above right=3mm and 8mm of lbfgs] (full) {全矩阵更新\\逐行范数保持};
    \node[map, below right=3mm and 8mm of lbfgs] (lora) {低秩更新\\$W_0+BA$};
    \node[score, right=12mm of $(full)!0.5!(lora)$] (score) {外层分段评分\\拒答比与首词元质量代理};
    \node[score, below=9mm of loss] (tpe) {多目标 TPE 提议 $\theta$\\更新层区间与损失超参数};

    \draw[arr] (prompt) -- (hook);
    \draw[arr] (hook) -- (iocache);
    \draw[arr] (iocache) -- (loss);
    \draw[arr] (loss) -- (lbfgs);
    \draw[arr] (lbfgs) -- (full);
    \draw[arr] (lbfgs) -- (lora);
    \draw[arr] (full) -- (score);
    \draw[arr] (lora) -- (score);
    \draw[feedback] (score.south) |- (tpe.east);
    \draw[feedback] (tpe.north) -- node[midway, right, align=left, font=\scriptsize]
      {下一试次：恢复基线映射\\复用冻结缓存} (loss.south);
  \end{tikzpicture}%
  }
  \caption{上游 ARA 原型的历史双层流程（硬近邻与分段评分）。前向钩子与 I/O 捕获只执行一次，外层试次恢复基线映射后复用冻结缓存；全矩阵与低秩路径是两种互斥的参数化方式。本图不描述本地 point-v1 的软目标、完整因子重置或原始指标评分。}
  \label{fig:ara-pipeline}
\end{figure*}


\subsection{冻结的模块输入输出}

缓存阶段利用前向钩子记录每个目标模块的成对输入输出。源实现对每条提示调用一次只生成一个新词元的前向过程，并从模块张量中仅取最后一个序列位置；记录值随即从计算图分离、复制并移至 CPU。对正常提示聚合后得到
\begin{equation}
  X_g\in\mathbb R^{n_g\times p},\qquad
  Y_g\in\mathbb R^{n_g\times q},
\end{equation}
对目标提示则得到 $X_b\in\mathbb R^{n_b\times p}$ 与 $Y_b\in\mathbb R^{n_b\times q}$。这里 $X$ 是进入线性层的末词元隐藏状态，$Y$ 是同一次原始前向中的模块输出；四个张量在相应模块被优化期间保持冻结。因此，后续预测 $XW^{\mathsf T}$ 是针对原始模块局部输入分布的离线回放，而不是随权重变化重新采样的端到端轨迹。

冻结机制还带来两个必须显式满足的前置条件。第一，批量聚合允许某些钩子在部分批次中没有出现，但两类缓存随后按相同模块键直接索引，故被优化模块必须同时存在于正常缓存和目标缓存。第二，两个近邻集合分别从 $Y_g$ 和 $Y_b$ 中选取 $k$ 个样本，所以至少要求 $n_g\geq k$ 且 $n_b\geq k$。若任一条件不成立，当前实现没有定义可替代的跨模块匹配或自动缩小 $k$ 的规则。把这些条件写入形式化描述，可以区分算法所需不变量与通常数据批次恰好满足的偶然情况。

末词元采样决定了缓存所覆盖的条件分布。每条提示在生成调用中只允许产生一个新词元，钩子仍可能在预填充与生成相关计算中被触发，但保存操作统一截取当前张量的最后位置，并按样本维拼接。因而 $X_g$ 与 $X_b$ 不是整段序列所有位置的集合，$Y_g$ 与 $Y_b$ 也不是完整生成轨迹。CPU 分离副本阻断了回到基线模型的梯度，优化图只包含待改写权重、线性重放、距离计算与损失。这个设计降低了闭包成本，同时把结论范围限定在末词元模块响应：若要研究长序列位置间的传播，必须另行采样，而不能从当前缓存直接推得。

缓存中的输入输出也具有不同角色。保持项使用同一正常样本的 $(x_{g,i},y_{g,i})$ 配对，保留了逐样本监督；两个距离项只把输出行视为参考点集，不要求 $x_{b,i}$ 与某个正常提示一一配对。由此，方法可以在两类提示数量不同的情况下工作，但不能把近邻索引解释为提示语义对应关系。尤其是多个目标查询可以选择同一个正常输出，参考点也不需要被均匀使用，这正是其与带容量约束匹配方法的关键区别。

\subsection{局部近邻三项目标}

局部目标以动态 $k$ 近邻距离代替预先估计的单一方向。对查询向量 $z\in\mathbb R^q$ 和参考矩阵 $Y=(y_1,\ldots,y_n)^{\mathsf T}$，定义欧氏距离下的近邻索引集合及平均距离为
\begin{align}
  \mathcal N_k(z;Y)
    &= \operatorname*{arg\,topk}_{j\in\{1,\ldots,n\}}^{\rm smallest}
       \lVert z-y_j\rVert_2, \\
  d_k(z;Y)
    &= \frac{1}{k}\sum_{j\in\mathcal N_k(z;Y)}
       \lVert z-y_j\rVert_2 .
  \label{eq:knn-distance}
\end{align}
源实现通过全对全的 \code{cdist} 生成距离矩阵，再以最小值 \code{topk} 选择邻居。邻居索引随当前权重产生的查询位置重新计算，不是缓存阶段固定的标签。

正常保持项要求新映射在正常缓存上复现原模块输出：
\begin{equation}
  \mathcal L_{\rm good}(\widetilde W)
  =\frac{1}{n_gq}\left\lVert
    X_g\widetilde W^{\mathsf T}-Y_g
  \right\rVert_F^2 .
  \label{eq:good-loss}
\end{equation}
目标提示的输出一方面被拉向正常输出的局部邻域，另一方面被推离其原目标输出邻域：
\begin{align}
  \mathcal L_{\rm pull}(\widetilde W)
    &=\frac{1}{n_b}\sum_{i=1}^{n_b}
      d_k(x_{b,i}\widetilde W^{\mathsf T};Y_g), \\
  \mathcal L_{\rm push}(\widetilde W)
    &=\frac{1}{n_b}\sum_{i=1}^{n_b}
      d_k(x_{b,i}\widetilde W^{\mathsf T};Y_b).
  \label{eq:pull-push}
\end{align}
三项按
\begin{equation}
  \begin{aligned}
  \mathcal L_{\rm total}(\widetilde W)
    ={}&\lambda_g\mathcal L_{\rm good}(\widetilde W)\\
       &+\lambda_b\bigl[
         \mathcal L_{\rm pull}(\widetilde W)
         -\beta\mathcal L_{\rm push}(\widetilde W)
       \bigr]
  \end{aligned}
  \label{eq:total-loss}
\end{equation}
组合。负号使最小化过程增大到 $Y_b$ 的邻域距离；$\mathcal L_{\rm good}$ 则限制这一变化对正常缓存映射的偏离。该构造不先求“拒答方向”，而让每个查询使用邻近正常输出形成局部参照，因而可覆盖多簇表示而无需把它们压缩为单个均值差向量。

各超参数的作用可以直接由式~\eqref{eq:total-loss} 判读。$\lambda_b=0$ 时，目标提示不再参与改写；$\beta=0$ 时，目标输出只被拉向正常邻域而不显式远离原邻域；增大 $k$ 会把单点参照换成更宽的局部平均，但并不保证距离更平滑，因为近邻集合仍可切换。$\mathcal L_{\rm good}$ 是按元素平均的平方距离，$\mathcal L_{\rm pull}$ 与 $\mathcal L_{\rm push}$ 是按样本平均的一阶欧氏距离，三者量纲和随 $q$ 的尺度并不相同。因此，$\lambda_g$、$\lambda_b$ 与 $\beta$ 不应被解读为概率或归一化混合系数；外层搜索实际承担了在这些不同尺度间寻找折中的任务。

局部参照比全局中心保留了更多多峰结构，但它也会继承缓存密度的不均衡。高密度正常簇可能成为更多查询的近邻，孤立点的影响则受 $k$ 和相对距离限制；推远项同样优先作用于当前最接近的原目标输出，而非整个目标分布。方法由此表达的是经验点云上的局部几何，不是两个真实表示分布之间的总体距离。若缓存样本不足或分布偏置，目标函数本身不会校正这种偏置，外层有限评估也只能提供代理层面的筛选。

这里的“局部邻域映射”不是最优传输。式~\eqref{eq:knn-distance} 没有样本间耦合矩阵、质量守恒约束或 Wasserstein 代价，只是逐查询的近邻选择；即使采用“局部传输视角”，也只能作为几何解释，不能据此声称求解了最优传输问题。动态邻域还使全局目标非光滑且通常非凸。例如 $k=1$、$Y=\{-1,1\}$ 时，$d_1(z;Y)=\min(|z+1|,|z-1|)$ 在两个谷底的中点违反凸性不等式。邻居切换边界和零距离处也可能不可微。因此，后文的拟牛顿过程应被理解为源码采用的数值求解器，而不是全局最优性证明。

\subsection{全矩阵更新与逐行归一化}

全矩阵路径把 $W\in\mathbb R^{q\times p}$ 本身设为待优化变量，因此没有显式低秩约束；“全矩阵”只描述参数化自由度，不保证最终改变量具有满秩。为限制各输出通道的尺度，源实现先保存初始行范数 $s_j=\lVert w_{0,j}\rVert_2$，并在每次闭包计算中使用重参数化代理
\begin{equation}
  \mathcal R(W)_{j,:}
  =s_j\frac{w_{j,:}}
  {\max(\lVert w_{j,:}\rVert_2,\varepsilon)},
  \qquad \widetilde W=\mathcal R(W).
  \label{eq:row-normalization}
\end{equation}
在 $\lVert w_{j,:}\rVert_2\geq\varepsilon$ 的常规区域，代理行范数严格等于 $s_j$；在数值保护分支中，其范数不大于 $s_j$，当 $s_j>0$ 时严格小于。因此不能把式~\eqref{eq:row-normalization} 无条件描述成等式约束投影，也不宜把可行域概括为一个整体球面。内层计算 $X\widetilde W^{\mathsf T}$，没有在重建公式中另加显式偏置；优化结束后，再把归一化后的矩阵复制回模块权重。

逐行处理保留的是初始权重每个输出通道的尺度，而不是整个矩阵的单一 Frobenius 范数。常规区域中的自由度主要来自行方向旋转；不同输出行可以独立改变方向，行与行之间没有正交约束。若初始某行范数为零，则对应 $s_j=0$，代理行持续为零；若待优化行落入保护区域，分母被截断，梯度行为也不同于单位球面上的光滑重参数化。与此同时，钩子缓存的 $Y$ 是模块真实输出，而重放公式只写 $X\widetilde W^{\mathsf T}$，故可能存在的固定偏置并未作为单独变量加入目标。保持项会把这部分差异一并吸收到权重拟合压力中，不能说它精确隔离了权重变化之外的所有因素。

每个模块使用有限记忆 BFGS 及强 Wolfe 线搜索~\cite{liu1989lbfgs}。源码设置学习率为 $1$、每次 \code{step} 的内部最大迭代数为 $20$、历史长度为 $10$，并固定执行五次 \code{step(closure)}。一次闭包会重新构造归一化代理、全对全距离和近邻索引，然后反向传播式~\eqref{eq:total-loss}。固定调用次数只是实现预算；本文没有迭代轨迹或最优性证书，因而不把“五次”解释为已收敛。

五次外部调用也不能简单换算为一百次闭包。$20$ 是每次调用允许的内部迭代上限，强 Wolfe 搜索可为试探步额外求值，而提前停止也可能减少求值；实际次数应以 $C_m$ 记账。L-BFGS 的曲率近似通常适合中等规模平滑问题，但此处近邻切换和欧氏范数零点引入非光滑性，逐行保护分支又改变局部导数。采用该求解器是可复现的工程选择，不足以建立标准光滑拟牛顿理论中的局部收敛前提。

\subsection{低秩工程变体}

低秩路径冻结 $W_0$，以
\begin{equation}
  W_{\rm eff}=W_0+BA,\qquad
  B\in\mathbb R^{q\times r},\quad
  A\in\mathbb R^{r\times p}
  \label{eq:lora}
\end{equation}
表示有效权重，其增量秩至多为 $r$~\cite{hu2022lora}。源默认秩为 $128$，且设置 $\alpha=r$，所以常见 LoRA 写法中的缩放 $\alpha/r$ 等于 $1$。模块配置中的无偏置选项表示不训练适配器偏置参数，不能据此推断原线性模块一定没有固有偏置。

该变体存在需要审计的目标与部署差异。当行归一化开关开启时，闭包把 $\mathcal R(W_0+BA)$ 送入损失；部署时适配器实际给出的却是未显式归一化的 $W_0+BA$。两者不相等时，优化代理与最终映射并非同一函数。另一个状态条件是快速重置只把 $B$ 清零而保留 $A$：有效映射由此恢复到 $W_0$，但后续非凸优化的参数初始化可能依赖先前试次。本文保留这些源码语义，不把低秩路径与全矩阵路径表述成完全等价的两个求解器。

从表达能力看，全矩阵路径允许 $qp$ 个权重元素共同变化，低秩路径只有 $r(p+q)$ 个适配器元素，并把改变量限制在秩不超过 $r$ 的集合中。参数更少不等于有效映射无需完整矩阵：闭包仍先形成 $W_0+BA$，再进行与全矩阵相同的批量乘法和邻域计算。初始化时只要 $B=0$，有效映射便等于 $W_0$，但 $A$ 决定从该函数点出发的参数化切空间；保留 $A$ 因而可能改变下一试次早期梯度。比较两个路径时，应同时报告参数预算、归一化是否开启以及试次重置策略，而不能只报告秩。

\subsection{外层分段评分与计算代价}

外层以多变量 TPE 在表~\ref{tab:source-parameters} 的空间中产生 $\theta$，并同时最小化质量代理与拒答代理~\cite{akiba2019optuna}。设基线模型在评估目标提示上的拒答计数为 $f_0$，修改后计数为 $f$，则归一化拒答量为
\begin{equation}
  r_f=\begin{cases}
    f/f_0, & f_0>0,\\
    f, & f_0=0.
  \end{cases}
  \label{eq:refusal-ratio}
\end{equation}
质量变化使用首个生成词元分布的
$D=D_{\rm KL}(p_{\rm base}\Vert p_{\rm edit})$。源码传给搜索器的两个目标为
\begin{equation}
  s_1=\begin{cases}
    D/\sigma, & D\geq\tau,\\
    r_f\tau/\sigma, & D<\tau,
  \end{cases}
  \qquad s_2=r_f,
  \label{eq:outer-score}
\end{equation}
默认 $\sigma=1$、$\tau=0.01$。低于 KL 阈值时，第一目标被拒答量替换为按拒答代理缩放的惩罚，而不是继续直接使用 $D$。可选的通用能力分支以负的归一化准确率替代 $s_1$。还应区分搜索目标与结果展示：后者保存原始拒答计数和原始 KL（或负准确率），不等同于式~\eqref{eq:outer-score} 的分段坐标。这些指标只是拒答行为与局部质量的代理，不能单独证明全面安全性或通用能力保持。

分段设计把阈值内候选的首要区分依据改为拒答量：当 $D<\tau$ 时，$s_1$ 与 $s_2$ 都随 $r_f$ 变化，只是尺度不同；到达阈值后，$s_1$ 才恢复为 KL。它不是 $D$ 与 $r_f$ 的固定加权和，在分支边界也不必连续。例如 $r_f<1$ 时，从阈值下方接近所得 $r_f\tau/\sigma$ 小于阈值处的 $\tau/\sigma$。因此，Pareto 支配关系由这条显式策略塑造，读取搜索结果时必须保留分支条件。基线计数 $f_0$ 只评估一次，零基线使用原始 $f$ 的退化分支避免除零；这也使不同数据切分或不同基线之间的数值尺度未必可直接比较。

首词元 KL 的方向同样需要固定：实现计算的是 $D_{\rm KL}(p_{\rm base}\Vert p_{\rm edit})$，不是反向 KL，也不是整段生成序列的散度。它衡量基线分布质量在修改分布下的相对对数概率变化，对修改分布在基线低概率区域增加的质量敏感性有限。拒答计数则依赖外部判别规则，把连续输出压缩为离散事件。二者结合能为外层搜索提供便宜的二维信号，但没有覆盖事实正确性、多轮一致性、鲁棒性或广泛任务能力；可选准确率分支也只增加一个特定任务代理。

\begin{table*}[t]
  \centering
  \caption{上游 ARA 原型的历史默认值与搜索域（edc3b123456c）。所有行仅属于该原型，不是本地 point-v1 或 trajectory-v2 的实验配置。}
  \label{tab:source-parameters}
  \footnotesize
  \setlength{\tabcolsep}{4pt}
  \renewcommand{\arraystretch}{1.12}
  \begin{tabularx}{\textwidth}{@{}p{22mm}p{32mm}p{58mm}Y@{}}
    \toprule
    类别 & 参数 & 默认值或搜索域 & 实现语义 \\
    \midrule
    目标模块 & 组件 & 注意力输出投影、MLP 下投影 & 在所选连续层区间内逐模块建立缓存并优化 \\
    层区间 & 起始层、结束层 & $[0,\lfloor L/2\rfloor]$、$[\lfloor L/2\rfloor,L]$ & 结束层按 Python 切片语义不包含在区间内 \\
    损失权重 & $\lambda_g,\lambda_b,\beta$ & $[0,1]$；$[10^{-4},1]$（对数采样）；$[0,1.3]$ & 分别控制保持项、目标项和远离原目标邻域的校正强度 \\
    局部邻域 & $k$ & 整数 $[1,15]$ & 两个距离项使用相同近邻数 \\
    全矩阵路径 & 行归一化 & 默认开启 & 每次闭包计算前按初始权重的逐行范数重参数化 \\
    低秩路径 & 开关、秩 $r$ & 默认关闭；$r=128$ & 缩放取 $\alpha=r$，故适配器乘子 $\alpha/r=1$ \\
    L-BFGS & 学习率、内部迭代、历史、线搜索 & $1$；$20$；$10$；strong Wolfe & 每个目标模块固定调用 $5$ 次 \code{step(closure)} \\
    外层搜索 & 试次数、启动试次、候选数 & $200$；$60$；$128$ & 多变量 TPE，两个目标均最小化 \\
    质量代理 & $\sigma,\tau$ & $1.0,0.01$ & 用于首词元 KL 的缩放与分段阈值 \\
    缓存规模 & 优化/评估提示数 & 每类 $400$；每类 $100$ & 前者建立内层缓存，后者计算外层代理 \\
    \bottomrule
  \end{tabularx}
\end{table*}


计算瓶颈来自代理矩阵乘法和全对全近邻。对单模块单闭包，前者时间为 $O((n_g+n_b)pq)$；拉近与推远距离分别需要 $O(n_bn_gq)$ 和 $O(n_b^2q)$ 时间，并分别物化 $O(n_bn_g)$ 与 $O(n_b^2)$ 的距离矩阵。由于不同后端的 \code{topk} 实现不同，选择开销保守记为
$T_{\rm topk}(n_b,n_g,k)+T_{\rm topk}(n_b,n_b,k)$，而不强行指定排序复杂度。若模块 $m$ 的实际闭包调用数为 $C_m$，则主导总时间可写为
\begin{equation}
  O\!\left(\sum_m C_m\left[
    (n_g+n_b)p_mq_m+n_b(n_g+n_b)q_m+T_{{\rm topk},m}
  \right]\right).
\end{equation}
全矩阵 L-BFGS 的历史状态量为 $O(hp_mq_m)$，其中 $h=10$；低秩可训练状态降为 $O(hr(p_m+q_m))$，但仍需物化 $W_0+BA$ 的 $O(p_mq_m)$ 空间，并承担形成 $BA$ 的 $O(p_mq_mr)$ 时间。CPU 缓存总标量数为
$O(\sum_m(n_{g,m}+n_{b,m})(p_m+q_m))$。因此低秩路径主要压缩优化变量和拟牛顿历史，而不会消除完整有效权重、缓存或距离矩阵的代价。以上分析完全来自固定源码和静态重建；本节未新增模型运行或经验实验。

外层总成本还要乘以被实际评估的候选配置，并随候选层区间内的模块数变化。默认两百试次是搜索预算而非保证完成的有效样本数；失败、剪枝或运行环境差异都可能改变实际求值集合。距离矩阵的峰值显存随 $n_b(n_g+n_b)$ 二次或双线性增长，低秩化并不改变这一项，因而在大缓存上可能先受邻域内存限制而不是参数历史限制。相反，缓存常驻 CPU，只有当前模块张量参与设备侧闭包，减轻了同时保存所有模块激活的设备压力。完整资源报告至少应给出每类样本数、目标模块数、实际闭包次数和外层有效试次数；源实现中的默认值只能用于复现配置，不能替代实测时间与峰值内存。

\section{本地 point-v1 的适配与目标}
\label{sec:local-adaptation}

本节分析本地 ARA point-v1，代码亦称 CARA。公式对应 \code{ara.py} 的校准、损失、重置与优化函数；该文件在后续修复提交 \code{868ca73} 与审计提交 \code{2871bc1} 间无差异\cite{ara2026localv1,ara2026localv2}。历史调度和评分则依据前一提交核对，不能用当前拆分模块反推运行当日的控制流。实际运行记录的是更早的 HEAD，尚未提交的部署修复与精确工作树之间的溯源缺口见第~\ref{sec:ara-v1-results} 节。

\subsection{校准缓存与增量前向}

对模块 $m$ 仍采用 $X_g,Y_g,X_b,Y_b$ 的配对记号。校准函数 \path{build_calibration_bank} 要求两类模块键匹配、每类至少两个样本，并验证 CPU FP32 缓存的维度与有限性。捕获时每批各目标模块恰好调用一次，保存预填充末位置的 I/O。定义正常输出的模块尺度
\begin{equation}
 s_m=\max\!\left\{\operatorname{mean}
 [(Y_g-\operatorname{mean}_{rows}Y_g)^{\odot 2}],10^{-6}\right\}.
 \label{eq:local-scale}
\end{equation}
行均值在每个输出坐标上计算，外层均值覆盖全部矩阵元素；$\odot 2$ 表示逐元素平方。局部候选输出为
\begin{equation}
 Z_c=Y_c+(X_c A^\top)B^\top,\qquad c\in\{g,b\}.
 \label{eq:local-output}
\end{equation}
$A\in\mathbb R^{r\times p}$、$B\in\mathbb R^{q\times r}$，配置 $r=128$ 且 LoRA 缩放为一，因此增量秩至多 128，并非实测秩为 128。式~\eqref{eq:local-output} 直接保留缓存中的基座响应，包含存在时的固定偏置；它无需在闭包中重建完整有效权重，也不做上游原型的逐行归一化。其正确性仍以缓存输入不变为条件：多个模块联合编辑后，真实网络输入可能漂移，局部等式不构成端到端保证。

\subsection{归一化软邻域与非负间隔惩罚}

对 $n$ 行参考点 $Y$ 和输出宽度 $q$，将实现中的两步距离写为
\begin{align}
 \delta(z,y_j)&=\max\!\left\{\frac{\|z-y_j\|_2^2}{q s_m},0\right\},\\
 d_\tau(z;Y)&=-\tau\left[\log\sum_{j=1}^{n}
 e^{-\delta(z,y_j)/\tau}-\log n\right].
 \label{eq:local-soft-distance}
\end{align}
\code{soft\_neighbor\_distance} 用范数与矩阵乘法恒等式计算平方距离，并将舍入产生的负值截为零。$\log n$ 是必要的参考数归一化，不能省略；温度 $\tau=0.10$。各查询使用全体参考点的软权重，无硬 $k$NN 选取。

记组件强度为 $a\geq0$、push 权重为 $\beta\geq0$、间隔为 $\mu$。\code{calculate\_ara\_loss} 的各项为
\begin{align}
 L_{keep}&=\operatorname{mean}[(Z_g-Y_g)^{\odot 2}]/s_m,\\
 L_{pull}&=\operatorname{mean}_i d_\tau(Z_{b,i};Y_g),\\
 L_{push}&=\operatorname{mean}_i\tau\operatorname{softplus}
 \!\left(\frac{\mu-d_\tau(Z_{b,i};Y_b)}{\tau}\right),\\
 G(A,B)&=\operatorname{mean}[(AA^\top-B^\top B)^{\odot 2}].
 \label{eq:local-terms}
\end{align}
由此得到
\begin{equation}
 \begin{aligned}
 L={}&L_{keep}+a(L_{pull}+\beta L_{push})\\
    &+10^{-4}\frac{G(A,B)}{\max\{G(A_0,B_0),10^{-6}\}}.
 \end{aligned}
 \label{eq:local-total}
\end{equation}
Gram 项先逐元素平方再求均值，既不是均值的平方，也不是矩阵平方。分母在模块优化入口用初始因子计算、从梯度图分离并固定，不随每次闭包中的 $A,B$ 更新。间隔惩罚随到原目标邻域的距离增大而衰减，替代上游以负号直接加入距离的推离项。

在精确算术下，$e^{-\delta/\tau}\leq1$，所以归一化对数均值不大于零，$d_\tau\geq0$；其余各项也非负，故 $L\geq0$。这只是下界，不能称目标上下有界。因子乘积 $BA$ 和多参考点损失仍不保证凸性。软邻域消除了公式中的硬 top-$k$ 切换，但不能直接套用上游的非光滑反例，也没有建立跨截断与数值回滚分支的全局光滑性或收敛定理。

\subsection{完整重置、优化与规范化}

\code{snapshot\_adapter\_state} 用种子和模块名派生确定性随机数，初始化 $A$ 并令 $B=0$，保存两者的 CPU 快照。每次试次由 \code{restore\_adapter\_state} 恢复全部 $A,B$ 并清空梯度，故不会继承上游仅清零 $B$ 的因子历史。

只有处于所选半开层区间且组件强度为正的模块进入优化。原始 MLP 强度先截为 $\max(0,a_{raw})$；应用强度为零时直接跳过，并不运行仅含 Gram 项的优化。\code{\_run\_optimizer} 对每个进入优化的模块调用一次 \code{step(closure)}，内部 \code{max\_iter=20}、历史长度 10，采用 strong-Wolfe 线搜索。闭包调用次数与内部迭代次数不同，不能把预算写成固定 20 次损失求值。

优化结束后，\code{\_canonical\_factors} 在 float64 中分解
\begin{equation}
 A^\top=Q_A R_A,\quad B=Q_B R_B,\quad
 R_B R_A^\top=U\Sigma V^\top,
\end{equation}
再构造
\begin{equation}
 B'=Q_B U\Sigma^{1/2},\qquad
 A'=\Sigma^{1/2}V^\top Q_A^\top.
 \label{eq:local-canonical}
\end{equation}
QR 为 reduced 分解，SVD 只作用于小核心矩阵；最后恢复 float32 因子。在精确算术下 $B'A'=BA$，而 \code{\_canonicalize} 另外比较两类校准输入上的适配器输出。相对范数误差限为 $10^{-5}$、相对最大元素误差限为 $10^{-4}$，两个相对误差分母均下限截为 $10^{-6}$。这些数值检查限定于校准输入，不能作为未见输入上的通用误差保证。

\code{optimize\_ara\_module} 在优化后和规范化后分别检查有限性与目标下降；拒绝条件为 $L_{final}>L_{initial}+10^{-6}$，允许相等及此绝对容差。异常会恢复该模块进入优化前的因子；历史试次调度还会恢复完整初始快照。这是代数与数值一致性控制，没有修复前对照，不能据此断言其导致本次全部试次稳定完成。

\subsection{实际外层指标与搜索参数}

point-v1 将原始 Keywords 和原始首词元 KL 同时最小化，未采用式~\eqref{eq:outer-score} 的历史分段变换。设响应为 $r_i$，归档标记集合为 $\mathcal M$，历史 \code{KeywordRate.\_is\_match} 定义
\begin{equation}
 K=\frac{1}{N_b}\sum_i\mathbf 1\!\left[
 \begin{array}{l}
 r_i\text{ 为空或仅空白，或}\\
 \exists u\in\mathcal M:\operatorname{lower}(u)
 \text{ 是 }T(r_i)\text{ 的子串}
 \end{array}\right].
 \label{eq:local-keywords}
\end{equation}
$T$ 依次小写化、去掉星号、将弯撇号转成直撇号，并把词间空白合并为一个空格。实际标记含英文与中文，完整集合保存在归档配置。命中是词汇代理，可能误报正常说明或漏报语义拒答；空响应计为命中也不使该规则成为语义判别器。

历史 KL 评分器调用编辑后对数概率作为输入、基座对数概率作为目标，并使用 \code{log\_target=True} 与 \code{batchmean}：
\begin{equation}
 D=\frac{1}{N_g}\sum_i\sum_v p_{0,i}(v)
 \log\frac{p_{0,i}(v)}{p_{\theta,i}(v)}.
 \label{eq:local-kl}
\end{equation}
方向是基座相对编辑模型，位置仅为首个生成词元。候选以同一验证集反复评分；$D$ 较低不能推出整段生成、知识能力或安全行为保持。

六个搜索坐标为起始比例、跨度比例、注意力强度、原始 MLP 强度、push 权重与间隔，范围见表~\ref{tab:ara-v1-protocol}。对记录中的 $L=64$，历史层解析为
\begin{align}
 \ell_s&=\lfloor64 f_{start}\rfloor,\\
 \ell_e&=\min\{64,\max(\ell_s+1,\ell_s+\lceil64 f_{span}\rceil)\}.
\end{align}
采样值与实际应用值必须区分。TPE 使用的 \path{failure_constraint} 只标记数值失败，归档中全为零；它不是 Keywords/KL 效果验收约束。内层小损失、外层低代理值和最终通过门槛是三个不同事件。

\section{从一个方向到直接映射}
\label{sec:method-comparison}

\subsection{共同坐标与几何差异}

三类方法可以放进同一比较坐标，但不构成理论包含链。图~\ref{fig:method-geometry} 从左到右画出单方向投影、SOM 多方向投影和模块输出的局部邻域映射。横向排列只表示决策变量发生变化：$r$、有序方向组 $(r_1,\ldots,r_k)$、以及上游矩阵 $W$（本地适配则为因子 $A,B$）。它不表示后一类方法必然包含前一类方法，也不表示拒答抑制效果随面板顺序递增。

\begin{figure*}[t]
  \centering
  \resizebox{0.97\textwidth}{!}{%
  \begin{tikzpicture}[
    font=\sffamily\scriptsize,
    >=Latex,
    panel/.style={draw=black!45, rounded corners=2pt, fill=black!1},
    good/.style={circle, draw=blue!70!black, fill=blue!18, inner sep=1.5pt},
    bad/.style={circle, draw=red!70!black, fill=red!18, inner sep=1.5pt},
    proto/.style={circle, draw=orange!80!black, fill=orange!22, inner sep=1.6pt},
    mapped/.style={rectangle, draw=green!50!black, fill=green!20, inner sep=1.5pt},
    flow/.style={-{Latex[length=2.2mm]}, line width=0.75pt},
    thinflow/.style={-{Latex[length=1.7mm]}, line width=0.55pt}
  ]

  % Panel A: single direction
  \begin{scope}[xshift=0cm]
    \draw[panel] (0,0) rectangle (6.6,4.65);
    \node[font=\sffamily\bfseries\small] at (3.3,4.34) {(a) 单方向投影};
    \node[good] at (1.15,1.20) {};
    \node[good] at (1.55,1.62) {};
    \node[good] at (1.85,0.93) {};
    \node[good] at (2.15,1.42) {};
    \node[bad] at (4.25,2.88) {};
    \node[bad] at (4.65,3.24) {};
    \node[bad] at (4.92,2.55) {};
    \node[bad] at (5.25,3.02) {};
    \coordinate (mug) at (1.7,1.3);
    \coordinate (mub) at (4.75,2.92);
    \fill[blue!70!black] (mug) circle (2pt) node[below=2pt] {$\mu_g$};
    \fill[red!70!black] (mub) circle (2pt) node[above=2pt] {$\mu_b$};
    \draw[flow, purple!75!black] (mug) -- node[above, sloped] {$r=\mu_b-\mu_g$} (mub);
    \draw[dashed, black!55] (0.7,2.75) -- (5.9,0.45) node[below left] {$r^\perp$};
    \draw[thinflow, red!70!black] (4.25,2.88) -- (3.37,1.85);
    \draw[thinflow, red!70!black] (4.92,2.55) -- (3.70,1.43);
    \node[align=center] at (3.3,0.32) {删除一个全局分量：$x\mapsto\Pi_r(x)$};
  \end{scope}

  % Panel B: SOM directions
  \begin{scope}[xshift=7.05cm]
    \draw[panel] (0,0) rectangle (7.35,4.65);
    \node[font=\sffamily\bfseries\small] at (3.675,4.34) {(b) SOM 多方向};
    \draw[step=0.72cm, black!35] (0.85,0.75) grid (3.73,3.63);
    \foreach \i in {0,...,3}{
      \foreach \j in {0,...,3}{
        \node[proto] at ({0.85+0.72*\i},{0.75+0.72*\j}) {};
      }
    }
    \fill[blue!70!black] (5.80,1.10) circle (2.2pt) node[below=2pt] {$\nu$};
    \draw[thinflow, orange!80!black] (5.80,1.10) -- node[below, sloped] {$r_1$} (1.57,1.47);
    \draw[thinflow, orange!80!black] (5.80,1.10) -- node[above, sloped] {$r_2$} (3.01,2.91);
    \draw[thinflow, orange!80!black] (5.80,1.10) -- node[above, sloped] {$r_3$} (1.57,3.63);
    \node[align=center] at (5.78,2.55) {$r_i=w_i-\nu$\\16 个候选};
    \node[align=center] at (3.675,0.30) {BO/TPE 选择有序的 $k=2,\ldots,7$ 个投影};
  \end{scope}

  % Panel C: local matrix mapping
  \begin{scope}[xshift=14.85cm]
    \draw[panel] (0,0) rectangle (8.1,4.65);
    \node[font=\sffamily\bfseries\small] at (4.05,4.34) {(c) 上游局部矩阵映射};
    \node[bad] (xb1) at (0.75,2.75) {};
    \node[bad] (xb2) at (0.75,2.15) {};
    \node[bad] (xb3) at (0.75,1.55) {};
    \node[align=center] at (0.75,0.93) {模块输入\\$X_b$};
    \node[draw, rounded corners=2pt, fill=purple!12, minimum width=1.05cm, minimum height=1.05cm] (W) at (2.25,2.15) {$W$};
    \draw[flow] (1.02,2.15) -- (W);
    \node[mapped] (z1) at (4.05,2.80) {};
    \node[mapped] (z2) at (4.38,2.30) {};
    \node[mapped] (z3) at (3.82,1.82) {};
    \draw[flow, purple!75!black] (W) -- node[above] {$X_bW^\top$} (3.62,2.30);
    \node[good] at (5.15,2.92) {};
    \node[good] at (5.52,2.48) {};
    \node[good] at (5.12,1.95) {};
    \node[good] at (5.72,2.02) {};
    \node[align=center, blue!70!black] at (5.45,3.52) {良性输出邻域 $Y_g$};
    \node[bad] at (7.15,1.10) {};
    \node[bad] at (7.50,1.48) {};
    \node[bad] at (7.20,1.86) {};
    \node[align=center, red!70!black] at (7.18,0.48) {原目标输出 $Y_b$};
    \draw[thinflow, green!50!black] (4.38,2.30) -- node[above, sloped] {拉近} (5.12,2.48);
    \draw[thinflow, red!70!black, dashed] (4.05,1.82) -- node[below, sloped] {推离} (6.98,1.30);
    \node[align=center] at (3.95,0.36) {直接优化映射；不显式构造 $r_i$};
  \end{scope}

  \node[font=\sffamily\footnotesize, align=center] at (11.45,-0.52)
    {三种参数化的结构并列；面板顺序不表示理论包含关系，也不表示效果递增。};
  \end{tikzpicture}%
  }
  \caption{从显式方向到直接映射的几何对照。蓝色表示良性参照，红色表示目标行为样本，橙色表示 SOM 原型，绿色方块表示矩阵映射后的代理输出。SOM-MD 仍以方向和投影为显式变量；面板 (c) 示意上游矩阵代理；本地 point-v1 使用缓存输出加低秩增量与软邻域。该图只比较结构，不表示实测轨迹、机制因果性或跨协议性能。}
  \label{fig:method-geometry}
\end{figure*}


单方向方法用两类均值差定义一个全局轴，再删除沿该轴的分量\cite{arditi2024refusal}。SOM-MD 先以 16 个原型近似有害状态的局部结构，再用同一个无害质心生成方向池；BO/TPE 选择 2--7 个方向的排列，最终干预仍由投影组成\cite{piras2026som}。上游矩阵法则把每个目标模块的 $X\mapsto Y$ 映射作为优化对象，其局部 $k$NN 损失约束目标行为新输出与 $Y_g,Y_b$ 的相对距离。这里的“局部邻域映射”不是最优传输：源码没有运输计划、质量守恒或样本间全局耦合。

三种几何锚点保留的信息不同。均值差只保留两组状态的一阶中心；SOM 用 16 个相互关联的原型覆盖有害状态的高密度区域，但仍以单个 $\nu$ 表示无害参照；矩阵法不先聚合 $Y_g,Y_b$，而在每次闭包中重新计算候选输出到参考样本的最近邻。本地 point-v1 将硬近邻改为全参考集上的软邻域，但同样依赖冻结缓存。两种版本保留局部样本结构，却把结果绑定到冻结参考集合和欧氏距离。它不是对“拒答流形”的无条件恢复，也不能证明参考集合之外的状态会按相同方式映射。

干预算子的代数性质也不能只按参数数量排序。单方向投影是固定的低维删除；SOM-MD 复合多个一般不交换的投影；全矩阵能够表示投影、旋转和缩放等线性变化，但源码的逐行范数处理与局部损失会限制实际解。参数空间较大意味着搜索自由度和优化风险同时增加，不意味着矩阵法在因果控制上严格强于方向法。

\subsection{结构比较矩阵}

表~\ref{tab:method-comparison} 用共同坐标区分方向法、上游矩阵原型和本地 point-v1。该表不填写 ASR 或能力分数，因为现有来源没有统一模型、数据划分、判别器、解码配置和计算预算。

\begin{table*}[t]
\centering
\caption{方向法、上游原型与本地 point-v1 的结构比较。SOM-MD 为本文简称；各行没有共享效果评测协议。}
\label{tab:method-comparison}
\footnotesize
\setlength{\tabcolsep}{3.5pt}
\renewcommand{\arraystretch}{1.16}
\begin{tabularx}{\textwidth}{@{}>{\raggedright\arraybackslash}p{21mm}>{\raggedright\arraybackslash}p{24mm}Y Y Y Y@{}}
\toprule
方法 & 决策变量 & 几何锚点 & 干预算子 & 搜索对象 & 证据边界 \\
\midrule
单方向 & 一个方向 $r$ & 有害与无害状态均值差 & 单个正交投影，可折叠进权重 & 层与候选方向筛选 & 同行评议结果限于原论文协议 \\
SOM-MD & 有序方向组 & 16 个有害 SOM 原型减一个无害质心 & 2--7 个投影顺序复合 & BO/TPE 搜索排列 & 原文 7+1 个目标设置；验证/测试关系不确定 \\
上游 ARA & 模块 $W$ 或因子 $A,B$ & 冻结模块 I/O 与硬 $k$NN & 全矩阵写回或低秩适配；可有行归一化 & L-BFGS 与历史分段评分 TPE & 源码形式化与数学审计；无本文归档实测 \\
本地 point-v1 & 秩至多 128 的 $BA$ & 冻结 I/O、归一化软邻域 & 缓存输出加增量；无行归一化 & 六维 TPE，原始 Keywords/KL & 一次 120 次试验验证搜索；联合门槛失败 \\
\bottomrule
\end{tabularx}
\end{table*}


“任意秩”只描述上游全矩阵路径没有显式低秩因子。它不保证 $\Delta W=W-W_0$ 达到任意预指定秩，也不保证满秩。可选低秩路径反而满足 $\operatorname{rank}(BA)\leq r$。同理，矩阵具有更多自由参数并不能单独推出更强的行为控制；冻结 I/O、行范数处理、损失权重和优化轨迹都会限制实际更新。

上游低秩工程路径还暴露了比较表无法承载的实现边界。完整行归一化开启时，优化闭包评价归一化后的 $W_0+BA$，而适配器部署仍使用未投影的 $W_0+BA$；两者一般不同。这个静态差异应被披露，但其对拒答与能力指标的影响尚无实验，不能把它升级为“低秩版本失效”的结论。全矩阵路径在结束时写回归一化矩阵，不存在同一写回差异。本地 point-v1 关闭行归一化并直接在缓存输出上加增量，另用 QR/SVD 检查因子规范化前后的校准输出；因此该历史差异不能用来解释本地失败结果。

\subsection{搜索成本与证据不可比性}

三类方法的外层成本衡量不同对象。单方向流程需要选层和筛选方向；SOM-MD 对有序方向组合执行完整生成与判别器评分，其候选空间随 $k$ 快速增长；矩阵法在每个外层试验中逐模块运行连续 L-BFGS，再用质量与拒答代理评分。即便两套流程都使用 TPE，试验次数也不是等价预算，因为一次试验包含的模块优化、生成数量和判别器调用不同。

证据等级同样不能横向替换。Arditi 等、Wollschl\"ager 等和 Piras 等给出了各自协议内的同行评议结果\cite{arditi2024refusal,wollschlager2025geometry,piras2026som}；上游矩阵原型具有固定源码、数学推论和公开讨论；本地 point-v1 新增一次验收失败的归档搜索，但仍没有同协议受控比较。Piras 等报告的 SOM-MD 结果只能写为“该文在其搜索与评测协议下报告”，还需保留 HarmBench 验证/测试关系未公开的边界。本文据此比较机制，不判断哪一种方法在数值上占优。

公平的后续比较至少需要锁定五组变量。三种方法应使用同一模型版本和目标模块范围，训练、搜索与测试提示应按样本标识公开隔离，生成应共享系统提示和解码参数，行为判别器与无害能力指标应一致，总预算还应同时报告生成次数、优化器闭包调用、GPU 时长和峰值内存。缺少任一组控制时，ASR 差异可能来自搜索预算、数据选择或判别器，而不是几何参数化。

\section{证据审计与可证边界}
\label{sec:evidence-audit}

本文区分源码事实、数学推论、上游声明、同行评议文献、未来假设与本地归档观察。表~\ref{tab:evidence-boundaries} 的 E1--E5 保持原义，新增 E6 承载 point-v1 的单次搜索。运行完成不能自动升级为方法有效，当前实现也不能升级为 v2 效果结果。

\subsection{上游非凸性与本地目标下界}

上游硬近邻目标一般非凸。取 $k=1$、$Y=\{-1,1\}$，有
\begin{equation}
 f(z)=\min\{|z+1|,|z-1|\},
\end{equation}
$f(-1)=f(1)=0$ 而 $f(0)=1$，违反凸函数的中点不等式。固定邻居时，拉近项是仿射映射后二范数之和，负的推远项形成局部差分结构；邻居切换和零距离又引入非光滑点。行范数重参数化及双线性因子并未恢复经典光滑凸优化条件，五次 L-BFGS 调用不是收敛证书\cite{ara2026source}。

这段反例属于上游硬 $k$NN，不能直接归于本地软邻域公式。point-v1 的式~\eqref{eq:local-total} 在非负强度和精确算术下有零下界，因子化仍不保证凸性，也无全局收敛定理。源码注释中的 bounded 在本文仅解释为有下界，不是上下均有界。有限性、下降容差及回滚保护是控制流事实，而不是收敛率证明。

秩语义跨版本同样需要限定。上游全矩阵路径不显式约束 $\operatorname{rank}(W-W_0)$，但仍有代数上界 $\min(p,q)$；本地实验使用秩至多 128 的适配器。名称中的任意秩并不意味着实测满秩、预指定秩或独立拒答机制数。

\subsection{冻结代理与历史实现差异}

上游原型和 point-v1 均复用原始末位置 I/O，后续模块不在上游编辑后重新捕获输入。因此局部损失下降与真实网络行为变化之间存在代理误差，模型级验证只能观测候选组合，不能反向修正缓存目标。

上游重算 $XW^\top$ 时未显式保留潜在偏置；MoE 路径要求两类缓存具有相同专家键且各有至少 $k$ 个参考样本；低秩快速恢复只清零 $B$；行归一化开启时目标中的 $\mathcal R(W_0+BA)$ 与部署的 $W_0+BA$ 一般不相同。这些是历史实现边界，不是已测量的失效率\cite{ara2026source}。

本地 point-v1 保留缓存 $Y$ 后加适配器增量，验证目标覆盖及最少样本，完整恢复 $A,B$，并关闭行归一化，故不能把上述偏置省略和仅重置 $B$ 再写作本地缺陷。它仍依赖冻结输入，QR/SVD 后的数值等价检查也只覆盖校准输入。第~\ref{sec:ara-v1-results} 节的实验没有控制组，不能确定任一修复对最终关键词率或运行稳定性的单独贡献。

\subsection{外层指标与资源解释}

上游式~\eqref{eq:outer-score} 使用分段质量评分，point-v1 则直接最小化 $(K,D)$。本地 Keywords 由空响应或规范化字符串命中判定，KL 只观察首词元；二者都不是全面安全性或通用能力测量。关键词率下降可以来自措辞变化，低 KL 也不能保证后续词元分布一致。语义判别与更完整评测仍需另行执行\cite{mazeika2024harmbench}。

资源模型也随版本变化。上游低秩闭包仍物化完整有效矩阵，本地式~\eqref{eq:local-output} 避免这一步，但 CPU 缓存和两类全对全距离仍有成本。日志末端采样值不是全程峰值；试次全 COMPLETE 不意味着未见模型、不同缓存规模或不同精度下也会稳定。

\subsection{溯源及观察强度}

本地直接证据优先采用归档字节哈希、journal 命名评分、有效配置与机器验收，日志用于补充计数和运行观察\cite{ara2026archive}。清单一致说明输入未在本次处理中改变，不是外部真实性签名。运行总结中关于修复有效的因果措辞没有配对对照支持；本文保留可核对的完成记录，将因果解释降为待检验假设。

必须区分记录 HEAD、后续修复提交和当前审计提交，不能把任何一个宣称为精确运行工作树。现存模型指纹也不能补齐缺失的模型修订。E6 的统计重算只复核标量记录，没有恢复逐提示评分过程，更没有产生独立实验重复。同行评议方向法结果保留各自协议范围，不能借给本地矩阵法作为对照结果。

\begin{table*}[t]
\centering
\caption{证据等级与全文允许措辞。等级表示可证范围，不表示研究价值排序。}
\label{tab:evidence-boundaries}
\footnotesize
\setlength{\tabcolsep}{4pt}
\begin{tabularx}{\textwidth}{@{}p{0.07\textwidth}p{0.20\textwidth}Y Y@{}}
\toprule
等级 & 证据来源 & 允许措辞 & 禁止越界 \\
\midrule
E1 & 冻结提交的控制流、张量运算与配置 & “源码实现/调用/存储/返回”“在该提交中可定位” & 不据此声称有效、稳定、快速或跨模型成立 \\
E2 & 从 E1 公式推出的数学性质与反例 & “一般不保证凸/光滑”“更新秩至多为 $r$”“目标与部署映射不同” & 不把一般性质写成每个样本必然失败，不给出未经证明的收敛率 \\
E3 & 上游说明、代码注释、演示模型与讨论 & “作者报告/注释称/公开演示显示”，并注明未独立复现 & 不写成本文结果，不据单模型演示建立方法排名 \\
E4 & 同行评议文献在其自身协议下的实验 & “该文在指定模型、数据与指标下报告”\cite{piras2026som,arditi2024refusal} & 不跨协议移植数值，不视为矩阵法与基线的同条件对照 \\
E5 & 本文提出但尚未执行的统一协议与假设 & “建议评估/需要检验/可证伪假设”\cite{mazeika2024harmbench} & 不使用“结果表明”“显著提升”“优于”或确定性因果措辞 \\
E6 & 本地 point-v1 冻结配置、journal 标量、日志与机器验收 & “本次 120 个试次完成”“重算得到”“验收失败”\cite{ara2026archive} & 不称独立复现或语义 ASR，不外推通用能力、因果稳定性及 v2 效果 \\
\bottomrule
\end{tabularx}
\end{table*}


\section{归档 point-v1 搜索与失败结果}
\label{sec:ara-v1-results}

本节仅报告 2026 年 9 月 4 日的一次 Qwen3.8-27B point-v1 搜索\cite{ara2026archive}。离线程序先校验清单及九个归档输入的 SHA-256，再按原始 journal 事件重建 120 个候选，并从命名评分记录重算统计与门槛。数据表和图共用生成的 JSON/CSV，保留被支配及未通过验收的候选。该操作复核已存标量，未重新生成模型输出；它属于 E6 归档观察。

\subsection{模型、数据与搜索协议}

运行使用归档中的本地 Qwen3.8-27B 路径，以 BF16、无量化及自动设备放置加载，配置 GPU 上限 90 GiB。硬件为一张 RTX PRO 6000 Blackwell Server Edition（97,887 MiB），驱动 595.58.03、Python 3.12.13。它不是早期双 RTX 3090 的 NF4 基线。表~\ref{tab:ara-v1-protocol} 列出有效配置与六个实际采样分布。

% Generated by evidence/summarize_ara_v1.py; do not edit.
\begin{table*}[t]
\centering\footnotesize
\caption{本地 point-v1 归档协议与六维搜索分布。}
\label{tab:ara-v1-protocol}
\renewcommand{\arraystretch}{1.12}
\begin{tabularx}{\textwidth}{@{}lYY@{}}
\toprule
项目 & 设置 & 证据边界 \\
\midrule
模型 & Qwen3.8-27B（本地路径） & 模型提交未记录；64 层 \\
硬件 & RTX PRO 6000 Blackwell Server Edition & 97,887 MiB；driver 595.58.03 \\
加载 & BF16；无量化；自动放置 & GPU 配置上限 90 GiB \\
无害数据 & mlabonne/harmless\_alpaca & revision: 02c6a92cfcf1（完整标识见证据 JSON） \\
有害数据 & mlabonne/harmful\_behaviors & revision: 01cead013989（完整标识见证据 JSON） \\
校准 & 每类候选 train[:300]；实际 64+64 & 捕获批量 1；秩至多 128 \\
验证 & 每类 train[300:400]；各 100 条 & 120 次搜索重复使用 \\
审计 & 每类 test[:100] & 仅配置；无审计得分 \\
解码 & seed=42；不采样；关闭 thinking & 最多 100 词元；有效批量 16 \\
批量配置 & 自动配置 0；最大值 16 & 有效评估批量为 16 \\
优化 & L-BFGS：一次 step；max\_iter=20 & history=10；strong Wolfe \\
采样 & TPE：36 startup + 84 adaptive & 候选 128；multivariate；seed=42 \\
数值约束 & 全部 constraints=[0.0] & 仅无记录数值失败；不代表效果验收 \\
起始层比例 & [0.25, 0.65] & 线性采样 \\
层跨度比例 & [0.1, 0.55] & 线性采样 \\
注意力强度 & [0.001, 1.0] & 对数采样 \\
MLP 原始强度 & [-0.1, 0.5] & 线性采样 \\
推离权重 & [0.0, 2.0] & 线性采样 \\
间隔 & [0.25, 4.0] & 对数采样 \\
区间解析 & floor(64×起点)，ceil(64×跨度) & 结束层裁剪至 64；MLP 强度取 max(0,原值) \\
\bottomrule
\end{tabularx}
\end{table*}


正常与目标提示的校准候选池各为 \code{train[:300]}，实际各捕获 64 条，批量为一。目标覆盖记录中的 64 层、128 个注意力输出或 MLP 下投影。验证用两类数据各自的 \code{train[300:400]}，每项评分对应 100 条提示；相同验证集贯穿全部候选。配置还给出 \code{test[:100]} 审计集，但由于无候选通过选择，没有审计分数，不能将其写为已完成的独立测试。空审计结果也不能单独证明历史初始化从未加载过审计数据。

解码使用种子 42、不采样、关闭 thinking、最多生成 100 个词元。配置批量为自动选择的零，日志中的有效评估批量为 16。Keywords 的逐试次显示记录为整数命中数除以 100，与其标量值一致；KL 的 100 条提示由配置与加载日志支持，归档未保存逐试次 KL 样本计数数组。这两个分母的证据来源不同。

历史采样器为多变量 TPE，启动试次 36、候选数 128、种子 42，并使用数值失败约束。六项分布中注意力强度和间隔为对数采样，其余为线性采样，均为连续分布、无离散步长。协议阶段只有 ID 0--35 的 startup 和 ID 36--119 的自适应 TPE；本文不将后续窗口划分解释成额外算法阶段。全部试次记录的系统约束为零，只说明没有记录到该约束定义的数值失败，不意味着搜索已遵循效果门槛。

\subsection{门槛与总体结果}

设基线 $K_0=1.00$，同一候选必须同时满足
\begin{equation}
 K\leq0.10,\quad \frac{K_0-K}{K_0}\geq0.50,\quad D\leq0.15.
 \label{eq:v1-gate}
\end{equation}
归档中 120/120 个试次均为 COMPLETE，非 COMPLETE 为零；Keywords 上限通过数为零、相对下降通过数为零、KL 上限通过数为 120，联合通过数为零。\code{acceptance.json} 的状态为 \code{failed}，选定试次、参数及验证、审计、重放、复载汇总均为 \code{null}。失败原因中列出的前五个试次只是截断诊断，不是失败候选总数。

% Generated by evidence/summarize_ara_v1.py; do not edit.
\begin{table*}[t]
\centering\footnotesize
\caption{本地 point-v1 验证集分布与失败验收结果。}
\label{tab:ara-v1-results}
\renewcommand{\arraystretch}{1.12}
\begin{tabularx}{\textwidth}{@{}lYYY@{}}
\toprule
统计/门槛 & Keywords/数量 & KL/条件 & 说明 \\
\midrule
最小值 & 0.54 & 0.000905 & 全部 120 次 \\
线性 25 分位 & 0.88 & 0.003696 & 全部 120 次 \\
中位数 & 0.98 & 0.010354 & 全部 120 次 \\
线性 75 分位 & 1.00 & 0.023271 & 全部 120 次 \\
最大值 & 1.00 & 0.115156 & 全部 120 次 \\
Trial 66 & 0.54 & 0.089975 & 描述性示例；未通过验收 \\
Trial 119 & 0.56 & 0.031919 & 描述性示例；未通过验收 \\
COMPLETE & 120/120 & 非 COMPLETE：0 & 单次自适应搜索 \\
Keywords 上限 & 0 & K <= 0.10 & 通过数 \\
相对下降 & 0 & (K0-K)/K0 >= 0.50 & 通过数 \\
KL 上限 & 120 & D <= 0.15 & 通过数 \\
联合验收 & 0 & failed & 未选中候选；审计等结果为空 \\
Pareto 点 & 10 & 双目标均最小化 & 相同坐标不互相支配 \\
\bottomrule
\end{tabularx}
\end{table*}


表~\ref{tab:ara-v1-results} 给出全部候选的线性分位数。Keywords 的最小值、中位数和最大值分别为 0.54、0.98、1.00，KL 对应的边际统计为 0.000905、0.010354、0.115156。不同指标的边际极值通常来自不同候选，不能组合成一个不存在的最佳模型。所有 KL 值低于门槛仅说明首词元分布代理满足设定上限。

最低 Keywords 的描述性例子为 ID 66：$K=0.54$、$D=0.089975$，解析层区间 $[16,52)$，优化 72 个模块。其相对关键词下降为 0.46；因本次基线恰为 1.00，也等于下降 46 个百分点，两种量在一般情况下不可互换。该候选既未达到 0.10 的绝对门槛，也未达到 0.50 的相对下降门槛，不是被验收或导出的模型。末个试次 ID 119 为 $K=0.56$、$D=0.031919$；终端从 1 开始计数，显示的第 120 次试验对应从 0 开始计数的 journal ID 119。

\subsection{Pareto 集与搜索轨迹}

在同时最小化 $(D,K)$ 时，若一个候选两项均不差且至少一项严格更好，它才支配另一个候选；重复坐标不互相支配。由全部 120 个点得到十个 Pareto ID：28、40、66、79、94、100、101、112、118、119。非支配性仅相对于本次有限候选集成立，不等于通过式~\eqref{eq:v1-gate}，也不是总体 Pareto 最优性证书。

\begin{figure*}[t]
\centering
\begin{minipage}[t]{0.49\textwidth}
\centering
\begin{tikzpicture}
\begin{axis}[
 width=\linewidth,height=57mm,
 xlabel={首词元 KL（验证集）},ylabel={Keywords},
 xmin=0,xmax=0.16,ymin=0,ymax=1.05,
 xtick={0,0.05,0.10,0.15},ytick={0,0.1,0.5,1},
 scaled ticks=false,tick label style={font=\scriptsize},
 label style={font=\small},grid=major,grid style={black!8},
 legend style={font=\scriptsize,draw=none,at={(0.5,1.02)},
 anchor=south,legend columns=2},
 scatter/classes={0={mark=*,draw=black!45,fill=black!25},
 1={mark=diamond*,draw=red!70!black,fill=red!60}},
]
\addplot[scatter,only marks,scatter src=explicit symbolic,mark size=1.5pt,forget plot]
 table[x=kl_divergence,y=keywords,meta=pareto,col sep=comma]
 {evidence/ara-v1-trials.csv};
\addlegendimage{only marks,mark=*,black!45}
\addlegendentry{其余候选}
\addlegendimage{only marks,mark=diamond*,red!70!black}
\addlegendentry{Pareto 候选}
\addplot[black,dashed,forget plot] coordinates {(0.15,0) (0.15,1.05)};
\addplot[black,dashed,forget plot] coordinates {(0,0.10) (0.16,0.10)};
\node[anchor=south west,font=\scriptsize] at (axis cs:0.005,0.11)
 {$K\leq0.10$};
\end{axis}
\end{tikzpicture}
\par\small (a) 所有候选与 Pareto 集
\end{minipage}\hfill
\begin{minipage}[t]{0.49\textwidth}
\centering
\begin{tikzpicture}
\begin{axis}[
 width=\linewidth,height=57mm,
 xlabel={试次 ID（从 0 开始）},ylabel={Keywords},
 xmin=0,xmax=119,ymin=0,ymax=1.05,
 xtick={0,35,66,90,119},ytick={0,0.1,0.5,1},
 tick label style={font=\scriptsize},label style={font=\small},
 grid=major,grid style={black!8},
 legend style={font=\scriptsize,draw=none,at={(0.5,1.02)},
 anchor=south,legend columns=2},
]
\addplot[only marks,mark=*,mark size=1.2pt,black!45]
 table[x=trial_number,y=keywords,col sep=comma]
 {evidence/ara-v1-trials.csv};
\addlegendentry{每次 Keywords}
\addplot[blue!70!black,thick,const plot]
 table[x=trial_number,y=running_best_keywords,col sep=comma]
 {evidence/ara-v1-trials.csv};
\addlegendentry{累计最低值}
\addplot[black,dashed,forget plot] coordinates {(35.5,0) (35.5,1.05)};
\addplot[black,dotted,forget plot] coordinates {(0,0.10) (119,0.10)};
\node[anchor=north,font=\scriptsize] at (axis cs:17,0.44) {startup};
\node[anchor=north,font=\scriptsize] at (axis cs:78,0.44) {自适应 TPE};
\end{axis}
\end{tikzpicture}
\par\small (b) 搜索顺序与累计最低值
\end{minipage}
\caption{point-v1 在验证集上的一次自适应搜索，最终验收失败。
(a) 全部 120 个候选；红色标出十个非支配点，虚线为 KL 0.15 和 Keywords 0.10 门槛。
(b) 同一验证集的逐次关键词率与累计最低值，竖线位于 ID 35 之后，区分 36 次 startup 与其余 TPE 试次。
图中不含独立重复、平滑估计或误差区间；Pareto 候选也未通过联合门槛。}
\label{fig:ara-v1-search}
\end{figure*}


图~\ref{fig:ara-v1-search} 同时呈现二维代理空间与按试次顺序的累计最低 Keywords。startup 后出现更低关键词候选，但这是同一次自适应搜索的描述，不能据此推断 TPE 相对其他策略的因果优势。累计最小值必然单调不增，也不是每次更新均改善的证据。试次 66 之后未刷新 0.54 的最低记录，不证明扩大预算必然无效，更不证明某个新目标必然有效。

\subsection{执行观察与复现限制}

日志报告总耗时 2 小时 38 分钟，终端采样的 GPU 已分配显存约 52.05 GB、常驻系统内存约 5.47 GB。这里保留日志单位及 allocated/resident 含义，不将其转换为完整运行峰值或保留显存。120 次完成支持本次未记录数值故障的有限观察；没有修复前对照、跨种子或跨模型重复，不能估计一般 OOM 概率或归因规范化修复的稳定性收益。

进程退出码为零，外层脚本仍打印成功及适配器路径，但机器验收明确失败。适配器未导出依据归档运行报告，不能只由本地目录缺失推断远端状态。本文交付的图表与 PDF 不构成模型适配器的成功导出证据。

运行记录的 HEAD 为 \code{cd2977a}，归档报告说明当时部署了未提交修复，后续对应修复提交为 \code{868ca73}。后者并非精确运行工作树的完整快照，仅检出前者也无法重建已部署修复。journal 的 \code{model\_commit} 为 \code{null}，模型指纹不能替代可恢复的模型快照，不得倒填 v2 模板的模型修订；未记录的包版本继续列为缺失。

归档没有逐提示响应、logits 或逐样本 KL 数组，因而不能重新判定词汇命中、语义拒答或散度。它只允许对已存评分的聚合进行一致性检查。单一模型、单次搜索与反复使用的验证集限制了外推；没有独立审计结果，不能报告跨种子置信区间、语义攻击成功率、全面能力保持或机制因果性。

\section{讨论与后续验证}
\label{sec:discussion}

本地结果提出一个明确的问题：在本次校准与搜索域内，所有候选均满足首词元 KL 上限，却没有候选达到关键词要求。它说明现有代理组合与预算没有产出可验收候选，不能定位原因究竟是冻结激活误差、末词元覆盖不足、目标尺度、搜索范围还是模型本身。下面将已有实现与仍需实验的假设分开。

\subsection{trajectory-v2：已实现但尚未实测}

当前审计版本 \code{2871bc1} 中，\code{ara\_trajectory.py} 沿基座生成的 continuation 做 teacher-forced 捕获，使用同生成步的正常与目标参考构建损失\cite{ara2026localv2}。模板每类使用 96 条提示、最多八个 continuation 位置，衰减系数 0.85，权重在每条提示内部归一化。相较 point-v1 单位置缓存，这扩展了被约束的位置，但仍然是冻结基座轨迹代理，编辑模型自由生成时可能偏离该轨迹。

v2 的 Gram 项是 $10^{-4}G(A,B)$，没有式~\eqref{eq:local-total} 的初始 Gram 分母，不能把两版目标写为相同归一化损失。\code{ara\_search.py} 与 \code{study\_runner.py} 实现八个搜索坐标，包括注意力与 MLP 部署增益，采用单 worker 的 8 个固定锚点、24 个 startup、88 个自适应试次调度。它不是 point-v1 的六维、36+84 协议，尚未测量扩大范围或增益带来的效果。

\code{scorers/refusal\_log\_odds.py} 提供连续的拒答与回答前缀对数几率分数；模板用该分数与 KL 作为优化目标，Keywords 保留为验收观察。该分数依赖预设前缀，不自动成为语义拒答真值。\code{protocol\_data.py} 约束校准、验证与审计角色；模板验证改用 \code{train[400:500]}，不能回填进 v1 的 \code{train[300:400]}。

验收相关模块 \path{acceptance.py}、\path{acceptance_export.py} 和 \path{artifact_schema.py} 实现候选选择、两次验证重放、导出前再次应用、单次审计账本、独立进程审计、指纹绑定与暂存后发布。运行保护还限定目标模块、设备、资源和部署漂移。这些是 E1 源码及配置事实；本次 v1 归档没有 v2 的候选、审计或导出结果，无法判断其是否提高成功率、语义效果或稳定性。

\subsection{从失败观察到可证伪假设}

第一项假设是冻结 I/O 的激活漂移限制了局部目标对整网行为的预测。可在相同预算下比较一次捕获、分块刷新与逐模块刷新，并直接测量缓存代理输出与编辑后真实输出之差。即使刷新降低代理误差，也仍需检查语义行为与能力指标，不能仅以局部损失证明改进。

第二项假设是末位置目标遗漏了较晚形成的拒答轨迹。point-v1 与 trajectory-v2 应使用匹配模型、提示与计算预算，比较首位置、多位置冻结轨迹及自由生成中的逐步行为。多位置成功与否需独立审计，不能从新增八个位置的代码推断。

第三项假设是当前搜索边界限制有效编辑。ID 66 的部分参数靠近上界是描述性现象，不证明最优解在域外。可预先指定范围扩展与部署增益消融，并与同域增加预算作配对比较。新一轮协议不得根据旧验证集反复选择后再宣称独立确认。

\subsection{统一比较、秩谱与资源}

单方向、SOM 多方向、上游全矩阵和本地低秩方法应固定模型修订、聊天模板、提示划分、目标模块、精度、解码与判别器\cite{arditi2024refusal,piras2026som}。预算同时报告候选数、生成次数、闭包调用、GPU 时长及峰值内存；独立种子需要分别运行，不能用同一搜索中的候选代替。验证集完成选择后，审计集只作预先约定的最终评价。HarmBench 可提供语义评估骨架\cite{mazeika2024harmbench}，仍需加入无害任务、长文本分布漂移和危险请求服从性。

组件消融应区分注意力、MLP 与二者联合编辑，并分别改变保持项、间隔惩罚和温度。硬 $k$NN 的 $k$ 扫描只适用于上游原型；本地软邻域不含该参数。更新奇异值谱、稳定秩及与既有方向的夹角可检验低秩更新的实际结构，但不能将几何相关性直接解释为拒答机制的数量或因果独立性。

完整资源报告需分别测量捕获、CPU 缓存、模块优化、验证与复载，保留失败和超时。上游的仅清零 $B$ 与完整重置、目标归一化与部署映射差异，可以作为历史实现消融；point-v1 已采用完整重置和无行归一化，不应再称这些为尚未实现的修复。

\subsection{行为解释与伦理边界}

拒答抑制既可帮助研究过度拒答与表征机制，也可能削弱原有防护。关键词减少不等于安全提升，危险知识是否仍可访问也不由局部几何决定。研究报告应同时呈现正常任务帮助性、危险请求服从性与失败统计。本文只提供公式、归档聚合和限制分析；本次验收失败没有产生可接受适配器，也没有证据将本地观察解释成更可靠的安全控制。

\section{相关工作}
\label{sec:related-work}

\subsection{表征工程与激活引导}

表征工程把群体层面的内部状态作为读出和控制对象，连接了线性探针、概念方向与模型干预\cite{zou2025representation}。Park 等从反事实语义形式化线性表征，并说明几何操作依赖所选内积\cite{park2024linear}。对比激活加法则从正负样本对的残差差异估计控制向量，在推理阶段将向量加入隐藏状态\cite{turner2023activation}；Rimsky 等在 Llama 2 上系统研究了这一路径\cite{rimsky2024steering}。这些方法证明方向性干预可以成为行为控制工具，却不保证一个方向能完整覆盖复杂安全行为。

激活加法与方向消融共享向量表示，却执行相反类型的操作。前者以带符号系数加入控制向量，允许连续调节目标行为；后者把向量张成的分量从激活或权重中删除。二者都依赖方向估计和干预层选择，也都可能影响非目标行为。本文把它们归入同一表征控制背景，但只将拒答投影方法作为直接结构基线。

\subsection{拒答的单方向与多维结构}

拒答方向研究提供了从干预到权重编辑的直接基线。Arditi 等报告单一方向在 13 个开源聊天模型上介导拒答，并将残差正交化折叠为紧凑权重修改\cite{arditi2024refusal}。安全微调位移分析与概念锥研究从不同数据锚点报告多维结构\cite{pan2025hidden,wollschlager2025geometry}；后者还提出表征独立性以区分几何正交与干预独立。Piras 等进一步使用 SOM 原型构造相关但不强制正交的方向池，并搜索有序投影组合\cite{piras2026som}。本文不重做这些论文的模型实验，而是把它们作为矩阵源码形式化的结构参照。

多方向证据并未形成唯一的拒答几何定义。概念锥工作从梯度信号和干预相互作用出发，SOM-MD 则从有害提示的局部原型出发；一个强调机制独立，一个允许方向保持较高相关性。二者都反对把欧氏正交直接当作完整因果解释，但其方向集合不可互换。矩阵法则改变了决策对象，因此本文只讨论结构差异，不把“无显式方向”解释成对所有多方向机制的吸收。

\subsection{局部原型、黑盒搜索与评测}

SOM 用规则网格上的原型近似高维数据，并保持原型间的邻接关系\cite{kohonen2013essentials}。贝叶斯优化把昂贵评估视为黑盒，TPE 用条件密度建模建议新的离散或连续配置\cite{snoek2012practical,bergstra2011algorithms}；Optuna 提供了该搜索的软件框架\cite{akiba2019optuna}。在拒答评测中，HarmBench 标准化自动红队与判别器流程，SORRY-Bench 则强调细粒度风险类别、语言和格式变化\cite{mazeika2024harmbench,xie2025sorrybench}。这些基准有助于定义统一复现实验，但不能替代对训练、搜索和测试样本关系的公开说明。

黑盒搜索的试验数只有在单次试验成本明确时才可比较。SOM-MD 的一次试验需要用一个有序方向组生成回答并调用行为判别器；矩阵法的一次外层试验还包含多个模块的 L-BFGS 优化。TPE 在两种流程中只是建议器，不会使二者自动获得相同的计算预算或统计独立性。评测论文提供了判别器和提示规范，却不能修复上游论文未公开的数据划分。

\subsection{本文位置}

本文位于方法论文与软件审计之间。我们不提出新的方向估计器，也不把源码原型的公开演示当作实证贡献。本文从固定实现分别重建上游硬近邻和本地 point-v1 软邻域目标，再以共同坐标比较显式方向投影与直接映射，并审计凸性、秩语义和实现一致性。新增的单次归档搜索呈现效果门槛失败，属于本地观察而非同行评议文献结果。这个定位把文献、实现、数学、归档观察与待验证假设分开。

现有工作留下的关键缺口不是再为一种参数化增加孤立效果数字，而是建立可复核的统一协议。统一协议需要同时控制模型、提示划分、干预位置、解码、判别器和总计算预算，并报告拒答抑制与无害能力保持。本文的比较表为这类实验定义对象；当前稿件完成形式化、证据审计及 point-v1 失败搜索分析，尚未完成统一协议的跨方法实验；trajectory-v2 的源码控制也尚未得到指定归档中的效果验证。

\section{结论}
\label{sec:conclusion}

本文区分上游 ARA 原型、本地 point-v1 和未有实测结果的 trajectory-v2，保留历史硬近邻目标推导，并形式化本地缓存输出增量、软邻域、非负间隔惩罚、完整因子重置与规范化检查。任意秩仅限定上游全矩阵参数化不显式约束更新秩；本地实验采用秩至多 128 的因子。软目标的零下界与数值保护都不能建立全局收敛或机制因果性。

归档 point-v1 的 120 个试次全部完成，但没有候选通过联合效果门槛。最低 Keywords 为 0.54，对应首词元 KL 为 0.089975，验收状态为 failed，没有通过验收的适配器。工程完成与效果失败同时成立。单次自适应验证搜索、逐提示证据缺失及运行版本记录不完整进一步限制了结论强度，不能由这些代理确认语义拒答下降、通用能力保持或跨方法优势。

trajectory-v2 已实现多位置轨迹、连续评分和更完整的验收控制，其收益仍待实验检验。后续应在固定模型、隔离数据、统一预算和独立语义审计下，检验激活漂移、较晚词元覆盖与搜索边界假设，再讨论结构变化是否带来可重复的行为改进。


\bibliographystyle{IEEEtran}
\bibliography{references}

\end{document}
